%% kapa: exam=1 round=24 year=2013 period=1 form=A subject=civil_law subject_ko=민법 seq=11 number=11
%% source: Q-Net 감정평가사 기출문제 (cst00309) artlSeq=5073802 file=extracted/1008/5073802/2013년 제24회 감정평가사 1차시험 기출문제/1교시(A형).hwp
%% conversion: kordoc 4.24.51 -> tools/gen1.py (auto); review=proofread-sonnet; flags=none
\begin{question}{11}
불공정한 법률행위에 관한 설명으로 옳지 않은 것은? (다툼이 있으면 판례에 의함)
\begin{choices}
  \item 경매에 의한 재산권의 이전에는 불공정한 법률행위의 법리가 적용되지 않는다.
  \item 대가관계 없이 당사자 일방이 상대방에게 일방적인 급부를 하는 법률행위도 불공정한 법률행위로 무효가 될 수 있다.
  \item 대리행위의 경우 경솔과 무경험은 대리인을 기준으로 판단하지만, 궁박은 본인을 기준으로 판단한다.
  \item 궁박은 경제적인 것에 한정되지 않으며, 정신적 또는 신체적 원인에 기인하는 것도 포함한다.
  \item 불공정한 법률행위가 성립되기 위한 요건은 무효를 주장하는 자가 증명해야 한다.
\end{choices}

eviewflag{choices 4-5 overlap in rhwp render; only part verified}
\end{question}
\answer{2}
